% Part: set-theory
% Chapter: ordinals
% Section: introduction

\documentclass[../../../include/open-logic-section]{subfiles}

\begin{document}

\olfileid{sth}{ordinals}{intro}

\olsection{مقدمه}

در~\olref[z][]{chap} با اصل~\stagesinf{} فرض کردیم که مرحله‌ای نامتناهی
از سلسله‌مراتب وجود دارد (اصل موضوع نامتناهی را نیز ببینید). اما با داشتن
\stagessucc{} نمی‌توانیم در مرحلهٔ نامتناهی متوقف شویم؛ باید ادامه دهیم.
پس در مرحلهٔ پس از نخستین مرحلهٔ نامتناهی، همهٔ گردایه‌های ممکن از
مجموعه‌هایی را می‌سازیم که در نخستین مرحلهٔ نامتناهی در دسترس بودند؛ و
دوباره؛ و دوباره؛ و دوباره؛ \dots

آنچه اینجا به‌طور ضمنی رخ داده این است که به استفاده از مفهومی «شهودی» از
عدد آغاز کرده‌ایم که بنا بر آن، ممکن است عددهایی \emph{پس از} همهٔ اعداد
طبیعی وجود داشته باشند. مفهوم مورد نظر، به‌ویژه، مفهوم \emph{عدد ترتیبی
ترامتناهی} است. هدف این فصل دقیق‌ترکردن این اندیشه است. مفهوم کلی عدد
ترتیبی را بررسی خواهیم کرد و سپس مجموعه‌های معینی را به‌صراحت به‌عنوان
اعداد ترتیبی خود تعریف خواهیم کرد.

\end{document}
